\section{Bitmask and DP}

\entry{PathDP (bitmask)}{$O(2^n n^2)$ \quad mem $2^n n$}{%
Folds every simple path covering \emph{exactly} a mask into one value, $n \le 20$.
Directed --- fill both directions for undirected. Pick a row of the table, edit the 5
config lines, done. The constructor does all the work; every call below is a
lookup.\\[0.3ex]
{\ttfamily\footnotesize
\begingroup\setlength{\tabcolsep}{2.5pt}
\begin{tabular}{@{}llllll@{}}
want & T & UNIT & NONE & ext(x,wt) & merge(a,b)\\
shortest path & ll & 0 & INF & x + wt & min(a, b)\\
longest path & ll & 0 & -INF & x + wt & max(a, b)\\
count paths & ll & 1 & 0 & x & a + b\\
count mod p & ll & 1 & 0 & x & (a + b) \% p\\
min bottleneck & ll & 0 & INF & max(x, wt) & min(a, b)\\
max probability & double & 1 & 0 & x * wt & max(a, b)\\
reachable lengths & bitset & 1<<0 & 0 & x << wt & a | b\\
\end{tabular}\endgroup}\\[0.3ex]
\begin{ctab}{0.25}{0.70}
\code{PathDP p(w,st,cn)} & \code{w[u][v]} = weight, \code{INF} = no edge.\\
\code{val(m,v)} & paths covering exactly \code{m}, ending at \code{v}.\\
\code{hamPath()} & over \code{v} of \code{val(full, v)}.\\
\code{hamPathTo(v)} & \code{val(full, v)}.\\
\code{hamCycle(s)} & closes \code{w[v][s]}; build with \code{st = 1LL<<s}.\\
\code{endingAt(v)} & over every mask, ending at \code{v}.\\
\code{bySize(k)} & over the masks of exactly \code{k} vertices.\\
\code{anyPath()} & over every mask and every \code{v}.\\
\code{cycles()} & over every simple cycle; build with \code{cn = true}.\\
\code{path(m,v)} & vertex order; \code{MIN} / \code{MAX} only.\\
\code{induced(g,vs)} & \code{w} for the vertices \code{vs}, weights all 1, local
\code{i} is \code{vs[i]}; translate back with \code{vs[x]}.\\
\end{ctab}\\[0.3ex]
\code{st} is a \textbf{mask}, not a vertex: \code{-1} = any start, \code{1LL<<s} pins
it. \code{cn} makes every path begin at the lowest vertex of its mask, for
\code{cycles()} only.\\[0.3ex]
\code{val(m,v)} is exactly \code{m}, not a subset, and \code{v} must be inside
\code{m}. Undirected counts come out $2\times$. \code{ext} must distribute over
\code{merge}, and \code{merge} must commute. Struct or bitset \code{T}:
\code{inline static const}, not \code{constexpr}. Unmodded counts overflow \code{ll} by
dense $n = 18$. At \code{-O2}, 30\% density: $n{=}16$ 0.01 s 13 MB, $n{=}18$ 0.05 s
44 MB, $n{=}20$ 0.22 s 184 MB.}

\begin{lstlisting}
struct PathDP {
    using T = ll;
    static T ext(T x, ll wt) { return x + wt; }
    static T merge(T a, T b) { return min(a, b); }
    static constexpr T UNIT = 0;
    static constexpr T NONE = 1e18;
    static constexpr ll INF = 1e18;      // missing edge in w
    int n; ll starts, full;
    matrix w; vector<T> dp; vector<char> par;
    PathDP(const matrix& wt, ll ss = -1, bool canon = false)
            : n(wt.size()), starts(ss), full((1LL << n) - 1),
              w(wt) {
        dp.assign((full + 1) * n, NONE);
        par.assign((full + 1) * n, -1);
        rep(v, 0, n) if (starts >> v & 1)
            dp[(1LL << v) * n + v] = UNIT;
        vector<ll> nb(n, 0);        // out-edges of u as bits
        rep(u, 0, n) rep(v, 0, n)
            if (w[u][v] < INF) nb[u] |= 1LL << v;
        rep(m, 1, full + 1) {
            ll base = m * n;
            for (ll b = m; b; b &= b - 1) {   // set bits of m
                int u = __builtin_ffsll(b) - 1;
                T cur = dp[base + u];
                if (cur == NONE) continue;
                for (ll c = nb[u] & ~m; c; c &= c - 1) {
                    int v = __builtin_ffsll(c) - 1;
                    if (canon && (1LL << v) < (m & -m))
                        continue;
                    ll t = (m | 1LL << v) * n + v;
                    T nv = merge(dp[t], ext(cur, w[u][v]));
                    if (nv != dp[t]) {
                        dp[t] = nv; par[t] = (char)u;
                    }
                }
            }
        }
    }
    T val(ll m, int v) { return dp[m * n + v]; }
    T hamPathTo(int v) { return val(full, v); }
    T hamPath() {
        T r = NONE;
        rep(v, 0, n) r = merge(r, val(full, v));
        return r;
    }
    T hamCycle(int s) {          // build with 1LL << s
        T r = NONE;
        rep(v, 0, n) {
            T d = val(full, v);
            if (d != NONE && w[v][s] < INF)
                r = merge(r, ext(d, w[v][s]));
        }
        return r;
    }
    T endingAt(int v) {
        T r = NONE;
        rep(m, 1, full + 1) r = merge(r, val(m, v));
        return r;
    }
    T bySize(int k) {            // exactly k vertices
        T r = NONE;
        rep(m, 1, full + 1) if (__builtin_popcountll(m) == k)
            rep(v, 0, n) r = merge(r, val(m, v));
        return r;
    }
    T anyPath() {
        T r = NONE;
        rep(m, 1, full + 1) rep(v, 0, n)
            r = merge(r, val(m, v));
        return r;
    }
    T cycles() {                 // build with canon
        T r = NONE;
        rep(m, 1, full + 1) {
            if (__builtin_popcountll(m) < 3) continue;
            int s = 0;
            while (!(m >> s & 1)) s++;   // lowest of m
            rep(v, 0, n) {
                T d = val(m, v);
                if (v != s && d != NONE && w[v][s] < INF)
                    r = merge(r, ext(d, w[v][s]));
            }
        }
        return r;
    }
    vi path(ll m, int v) {               // MIN / MAX only
        vi p;
        while (v >= 0) {
            p.push_back(v);
            int u = par[m * n + v]; m ^= 1LL << v; v = u;
        }
        reverse(p.begin(), p.end());
        return p;
    }
    template<class G, class V>    // g, vs: any integer type
    static matrix induced(const G& g, const V& vs) {
        static vi idx;        // reused, all -1 on entry
        if (idx.size() < g.size()) idx.assign(g.size(), -1);
        int k = (int)vs.size();
        rep(i, 0, k) idx[vs[i]] = i;
        matrix w(k, vi(k, INF));
        rep(i, 0, k) for (ll u : g[vs[i]])
            if (idx[u] >= 0) w[i][idx[u]] = 1;
        rep(i, 0, k) idx[vs[i]] = -1;   // wipe what we set
        return w;
    }
};
\end{lstlisting}

\entry{Subset sums (SOS / zeta)}{$O(2^n n)$ \quad mem $2^n$}{%
\code{sos(n, a)[m]} is the sum of \code{a[s]} over \emph{all} submasks \code{s} of
\code{m}. Replace \code{+=} with \code{|=} / \code{max} for ``or over submasks'' /
``best submask''. Supersets instead: flip the condition to
\code{if (!(m >> b \& 1)) f[m] += f[m | 1LL<<b]}.\\[0.3ex]
Inverse (M\"obius, recovers \code{a} from \code{f}): same loops with \code{-=}.
Subset convolution and ``count pairs with \code{a \& b == 0}'' both reduce to
this.\\[0.3ex]
To walk the submasks of one mask only: \code{for (ll s = m; s; s = (s-1) \& m)}, which
visits every nonempty submask once --- $O(3^n)$ over all masks.}

\begin{lstlisting}
// f has size 1<<n and holds a on entry
vi sos(int n, vi f) {
    rep(b, 0, n) rep(m, 0, 1LL << n)
        if (m >> b & 1) f[m] += f[m ^ (1LL << b)];
    return f;
}
\end{lstlisting}

\entry{LIS}{$O(n \log n)$ \quad mem $n$}{%
\code{c[k]} ends up as the smallest possible tail of an increasing subsequence of
length \code{k+1}, so \code{c.size()} is the answer. \code{c} itself is \emph{not} a
valid subsequence.\\[0.3ex]
\code{lower\_bound} gives \textbf{strictly} increasing; \code{upper\_bound} gives
non-decreasing. Longest decreasing: reverse the input, or negate it. Minimum number of
non-increasing chains covering the sequence = length of the LIS (Dilworth).\\[0.3ex]
To reconstruct, store for each element the index it was placed at plus the predecessor
\code{c}-index, then walk back from the last placement.}

\begin{lstlisting}
ll lis(const vi& a) {
    vi c;
    for (ll e : a) {
        auto it = lower_bound(all(c), e);
        if (it == c.end()) c.push_back(e);
        else *it = e;
    }
    return (ll)c.size();
}
\end{lstlisting}

\entry{Interval DP}{$O(n^3)$, optimised $O(n^2)$ \quad mem $n^2$}{%
Shortest interval first, so both halves are ready when you need them. Circular
versions: double the array to length $2n$ and take the best window of length
$n$.\\[0.3ex]
\textbf{Knuth} cuts it to $O(n^2)$ when \code{opt[i][j-1] <= opt[i][j] <= opt[i+1][j]}
--- restrict \code{k} to that window and keep an \code{opt} table. It holds when the
cost satisfies the quadrangle inequality, which is true for the classic ``merge
adjacent piles, cost = sum of the pile'' problem. \textbf{Divide and conquer
optimisation} is the same idea one layer at a time, for
\code{dp[i][j] = min over k of dp[i-1][k] + C(k,j)} with monotone \code{opt}:
recurse on \code{(jlo,jhi,klo,khi)}, evaluate the midpoint, split. Both are worth it
only when $n^3$ is too slow --- check the bound first.}

%NOCOMPILE
\begin{lstlisting}
rep(len, 2, n + 1)
  rep(i, 0, n - len + 1) {
    ll j = i + len - 1;                  // inclusive
    dp[i][j] = INF;
    rep(k, i, j)                         // split point
      dp[i][j] = min(dp[i][j],
                     dp[i][k] + dp[k+1][j] + cost(i, j));
  }
\end{lstlisting}

\entry{Divide and conquer DP}{KACTL \quad $O(n \log n)$ per layer}{%
The optimisation described above, as code. Shape: split \code{1..n} into \code{G}
consecutive groups, \code{cur[i] = min over k < i of prv[k] + cost(k, i)} where
\code{cost(k, i)} is the price of the group \code{(k, i]}, one layer per group. Naively
$O(Gn^2)$; this is $O(Gn\log n)$.\\[0.3ex]
Valid when the best \code{k} never moves left as \code{i} grows --- true when
\code{cost} satisfies the quadrangle inequality, e.g.\ (sum of the group)$^2$, or the
number of equal pairs inside it. Not sure? Compare with the naive DP on random small
inputs.\\[0.3ex]
Per layer: \code{prv} = previous layer (\code{prv[0] = 0}, the rest \code{INF} for
layer 1), \code{cur.assign(n+1, INF); DP().solve(1, n+1); prv = cur;}. \code{lo},
\code{hi} bound \code{k} to \code{[lo, hi)}. \code{cost} must be $O(1)$: prefix
sums.}

%LABEL dncdp
\begin{lstlisting}
vi prv, cur;          // previous layer, layer being built
ll cost(int k, int i);              // cost of group (k, i]
struct DP { // Modify at will:
    int lo(int ind) { return 0; }
    int hi(int ind) { return ind; }
    ll f(int ind, int k) { return prv[k] + cost(k, ind); }
    void store(int ind, int k, ll v) { cur[ind] = v; }
    void rec(int L, int R, int LO, int HI) {
        if (L >= R) return;
        int mid = (L + R) >> 1;
        pair<ll, ll> best(LLONG_MAX, LO);
        rep(k, max(LO,lo(mid)), min(HI,hi(mid)))
            best = min(best, make_pair(f(mid, k), k));
        store(mid, best.second, best.first);
        rec(L, mid, LO, best.second+1);
        rec(mid+1, R, best.second, HI);
    }
    void solve(int L, int R) { rec(L, R, INT_MIN, INT_MAX); }
};
\end{lstlisting}
